\documentclass[11pt]{amsart}
\usepackage{amssymb}
\usepackage[margin=1.15in]{geometry}
\usepackage[hidelinks]{hyperref}
\newtheorem{theorem}{Theorem}
\newtheorem{lemma}[theorem]{Lemma}
\theoremstyle{definition}
\newtheorem{definition}[theorem]{Definition}
\theoremstyle{remark}
\newtheorem{remark}[theorem]{Remark}
\newcommand{\code}[1]{\texttt{#1}}
\newcommand{\Gr}{\mathrm{Gr}}
\newcommand{\adj}{\mathrm{adj}}
\title{Sharpe's theorem on linear equations, and its Unit formalisation}
\date{October 26, 2026}
% stamp on every page: file name and date/time of compilation
\newcount\stamphour \newcount\stampmin
\stamphour=\time \divide\stamphour by 60
\stampmin=\time \multiply\stamphour by 60 \advance\stampmin by -\stamphour \divide\stamphour by 60
\edef\stamptime{\the\year-\ifnum\month<10 0\fi\the\month-\ifnum\day<10 0\fi\the\day\space
 \ifnum\stamphour<10 0\fi\the\stamphour:\ifnum\stampmin<10 0\fi\the\stampmin}
\AddToHook{shipout/foreground}{\put(0.5\paperwidth,-\dimexpr\paperheight-0.45in\relax)%
 {\makebox[0pt]{\scriptsize\ttfamily\jobname.pdf\quad\stamptime}}}
\begin{document}
\maketitle

Let $R$ be a commutative ring and $F$ an $R$-module.  A system of $m$
linear equations in $n$ unknowns
\[
  \textstyle\sum_{j=1}^n u_{ij}\,x_j=v_i\qquad(1\leq i\leq m),
\]
with $u_{ij}\in R$ and $v_i\in F$, to be solved in $F$, has at most one
solution when the ideal $D_n(U)$ of $n$-minors of $U=(u_{ij})$ has grade
at least~$1$ on $F$, and every solution makes the $(n+1)$-minors of
$W=[U\mid v]$ vanish.  Sharpe~\cite{Sharpe} proved that grade at
least~$2$ gives the converse: if $\Gr(D_n(U);F)\geq2$ and the
$(n+1)$-minors of $W$ vanish, the system has a solution.  The hypothesis
cannot be weakened to grade~$1$: over $\mathbb Z$ the system $2x=3$,
$4x=6$ has vanishing $2$-minors and $\Gr(2,4)\geq1$, but no solution.  The
theorem is Theorem B-7.5 of~\cite{CL}, where it is derived from a
local-global principle for coverings in depth~$2$.

This note gives a short direct proof, in which the depth-$2$ hypothesis is
used exactly once, on the Cramer numerators of the square subsystems, and
describes its checked formalisation in Unit (\code{sag/depth2}).  No
localisation, prime ideal or Noetherian hypothesis is used.

\section{Statement}

As in~\cite{CL,CQ}, depth is used in an explicit form, which for a
finitely generated ideal is equivalent to the true grade of Hochster and
Northcott~\cite{Northcott}.

\begin{definition}
Let $(d_\rho)_{\rho\in P}$ be a finite family of elements of $R$.
\begin{enumerate}
\item $\Gr(d;F)\geq1$ if $d_\rho x=0$ for all $\rho$ implies $x=0$
  ($x\in F$).
\item $\Gr(d;F)\geq2$ if moreover every family $(y_\rho)$ of elements of
  $F$ with $d_\rho y_{\rho'}=d_{\rho'}y_\rho$ for all $\rho,\rho'$ is of
  the form $y_\rho=d_\rho e$ for some $e\in F$.
\end{enumerate}
\end{definition}

Under $\Gr(d;F)\geq1$ the element $e$ in~(2) is unique.  We index the
minors by lists of rows.  For $\rho:\{1,\dots,n\}\to\{1,\dots,m\}$, let
$U_\rho$ be the $n\times n$ matrix with rows $\rho(1),\dots,\rho(n)$ of $U$,
and $\delta_\rho=\det U_\rho$.  Rows may repeat and appear in any order;
this gives $0$ or $\pm$ a minor, so the family $(\delta_\rho)$ generates
$D_n(U)$.  For $\sigma:\{1,\dots,n+1\}\to\{1,\dots,m\}$ and
$1\leq t\leq n+1$, let $c_\sigma(t)$ be the determinant of the
$(n+1)\times(n+1)$ matrix whose first $n$ columns are those of $U$ at the
rows $\sigma$ and whose last column is the unit vector $e_t$.  The
$(n+1)$-minor of $W$ at $\sigma$ is its expansion along the last column,
\[
  \textstyle\mu_\sigma=\sum_{t=1}^{n+1}c_\sigma(t)\,v_{\sigma(t)}\in F;
\]
for $F=R$ this is the determinant of the rows $\sigma$ of $W$.

\begin{theorem}[Sharpe]\label{thm:sharpe}
If $\Gr(\delta;F)\geq2$ and $\mu_\sigma=0$ for all $\sigma$, the system
$Ux=v$ has a unique solution in $F^n$.  Conversely, if it has a solution,
then $\mu_\sigma=0$ for all $\sigma$.
\end{theorem}

\section{Two determinant identities}

For a square matrix $A$ of size $n$, $A[k\leftarrow b]$ is $A$ with column
$k$ replaced by $b$, and $\adj(A)_{kt}=\det A[k\leftarrow e_t]$.  By
linearity of $\det$ in one column, $(\adj(A)\,b)_k=\det A[k\leftarrow b]$.

\begin{lemma}\label{lem:cramer}
$\adj(A)A=\det A\cdot I$ and $A\adj(A)=\det A\cdot I$.  Hence, for every
vector $b$, $\sum_k\det A[k\leftarrow b]\,a_k=\det A\cdot b$, where $a_k$
is the $k$-th column of $A$.
\end{lemma}

\begin{proof}
The first identity is column linearity together with the vanishing of a
determinant with a repeated column.  For the second, it suffices to show
$\adj(A^T)=\adj(A)^T$, i.e.\ $\det A^T[k\leftarrow e_i]=\det A[i\leftarrow
e_k]$.  Let $Z$ be $A$ with column $i$ replaced by $e_k$ and the other
entries of row $k$ replaced by $0$.  Then $A[i\leftarrow e_k]E=Z$ for the
elementary matrix $E=I+e_iw^T$ with $w_l=-A_{kl}$ for $l\neq i$ and
$w_i=0$, so $\det A[i\leftarrow e_k]=\det Z$.  The same construction for
$A^T$, $k$ and $i$ gives $Z^T$, so $\det A^T[k\leftarrow e_i]=\det Z^T=\det
Z$.
\end{proof}

\begin{lemma}\label{lem:rows}
Let $\rho$ be a list of $n$ rows, $r$ a row and $\sigma=(\rho,r)$.  For
every $v\in R^{n+1}$, writing $v'=(v_1,\dots,v_n)$ and $N_v$ for the
matrix with the columns of $U$ at the rows $\sigma$ and last column $v$,
\[
  \textstyle\det N_v+\sum_{j=1}^n u_{rj}\det U_\rho[j\leftarrow v']
  =\delta_\rho\,v_{n+1}.
\]
\end{lemma}

\begin{proof}
Let $N=N_{e_{n+1}}$.  The last row of $N$ is $(u_{r1},\dots,u_{rn},1)$
and its last column is $e_{n+1}$, so $\det N=\delta_\rho$.  Likewise
$N[j\leftarrow v]$ has last column $e_{n+1}$ and top left block
$U_\rho[j\leftarrow v']$, so its determinant is $\det U_\rho[j\leftarrow
v']$; and $N[n+1\leftarrow v]=N_v$.  The last coordinate of the identity
of Lemma~\ref{lem:cramer} for $N$ and $b=v$ is the claim.
\end{proof}

\section{Proof of the theorem}

Let $X_{\rho,k}=\sum_t\adj(U_\rho)_{kt}\,v_{\rho(t)}\in F$, Cramer's
numerators for the square subsystem of rows $\rho$, and write
$(Ux)_r=\sum_ju_{rj}x_j$ for $x\in F^n$.

\emph{The rows.}  Applying Lemma~\ref{lem:rows} to $v=e_t$, multiplying
by $v_{\sigma(t)}$ and summing over $t$ gives, for all $\rho$ and $r$,
\[
  \mu_{(\rho,r)}+(UX_\rho)_r=\delta_\rho\,v_r .\tag{1}
\]
On the left, the term $t=n+1$ of the second sum vanishes, since
$e_{n+1}'=0$; the terms $t\leq n$ give $X_\rho$.

\emph{Uniqueness.}  For $x\in F^n$,
$\sum_t\adj(U_\rho)_{kt}(Ux)_{\rho(t)}=(\adj(U_\rho)U_\rho\,x)_k
=\delta_\rho x_k$ by Lemma~\ref{lem:cramer}.  So if $Ux=Uy$, then
$\delta_\rho x_k=\delta_\rho y_k$ for all $\rho$, and $x=y$ as soon as
$\Gr(\delta;F)\geq1$.

\emph{Existence.}  Assume all $\mu_\sigma$ vanish, so that
$UX_\rho=\delta_\rho v$ for all $\rho$ by~(1).  Then
$U(\delta_{\rho'}X_\rho)=\delta_{\rho'}\delta_\rho v=U(\delta_\rho
X_{\rho'})$, and by uniqueness $\delta_{\rho'}X_{\rho,k}=\delta_\rho
X_{\rho',k}$ for every $k$.  So for each $k$ the family
$(X_{\rho,k})_\rho$ is proportional to $(\delta_\rho)_\rho$, and
$\Gr(\delta;F)\geq2$ gives $x_k\in F$ with $X_{\rho,k}=\delta_\rho x_k$
for all $\rho$.  Then
$\delta_\rho(Ux)_r=(UX_\rho)_r=\delta_\rho v_r$ for all $\rho$, hence
$Ux=v$.

\emph{Converse.}  If $Ux=v$, then
$\mu_\sigma=\sum_jx_j\sum_tc_\sigma(t)u_{\sigma(t)j}$, and the inner sum
is the determinant of the rows $\sigma$ of $W$ with its last column
replaced by column $j$ of $U$: a repeated column, so $0$.\qed

\begin{remark}
In~\cite{CL} the existence is obtained by solving the system over each
$R[1/\delta_\rho]$ by Cramer's rule and gluing the local solutions with
the local-global principle B-7.4 for coverings in depth~$2$; Sharpe first
treats an explicit regular sequence of length~$2$ in $D_n(U)$ and then
passes to a polynomial extension.  Here the local solutions are replaced
by the numerators $X_\rho$, whose compatibility is an identity
(uniqueness), and depth~$2$ is applied to them directly.  The
determinant identities that Sharpe establishes are Lemmas~\ref{lem:cramer}
and~\ref{lem:rows}.
\end{remark}

\section{The formalisation}

The theorem is checked in Unit, in \code{sag/depth2}, with univalence
and replacement as parameters as in the rest of \code{sag}
(\code{./segc sag/depth2 all}: 171 modules, about 2.6\,s; every file of
the development under 0.05\,s).  Matrices are lists of columns, as in
\code{sag/exterior}, where $\det$, $\adj(A)A=\det A\cdot I$,
$\det A^T=\det A$ and $\det(AB)=\det A\det B$ are proved.  Depth is the
explicit form above (\code{Depth}); for a family indexed by
$\{1,\dots,N\}$ it coincides with the Koszul grade of \code{sag/koszul}
(\code{DepthKoszul}).  The ring is the module \code{Modules.regular}, so
the case $F=R$ is an instance; it is also proved separately.

\medskip
\begin{tabular}{p{0.24\textwidth}p{0.68\textwidth}}
\code{Depth} & $\Gr\geq1$, $\Gr\geq2$, uniqueness of the multiple \\
\code{Cofactor} & Lemma~\ref{lem:cramer} (\code{symmetry}, \code{adj-right}, \code{cramer}), via \code{WiebeMatrices.emat} \\
\code{SharpeRows} & Lemma~\ref{lem:rows} (\code{Row}), via \code{WiebeMatrices.BlockDet} \\
\code{Sharpe} & Theorem~\ref{thm:sharpe} for $F=R$ (\code{System.theorem}, \code{sol-prop}, \code{necessary}) \\
\code{SharpeModule} & finite sums in a module, (1) (\code{RowF}), Theorem~\ref{thm:sharpe} (\code{System.theorem}, \code{sol-prop}, \code{necessary})
\end{tabular}
\medskip

The same directory contains the other results of~\cite[B-7]{CL} in
depth~$2$: the local-global principles for divisibility and strong gcds
(B-7.1), for depth~$2$ itself (B-7.2 with $k=2$), and the covering
principle (B-7.4), for localisations at single elements.

\begin{thebibliography}{9}
\bibitem{CL} T.~Coquand, H.~Lombardi, \emph{R\'esolutions libres finies.
  M\'ethodes constructives}, Calvage \& Mounet, 2024.
\bibitem{CQ} T.~Coquand, C.~Quitt\'e, \emph{Constructive finite free
  resolutions}, manuscripta math.~137 (2012), 331--345.
\bibitem{Northcott} D.~G.~Northcott, \emph{Finite Free Resolutions},
  Cambridge Tracts in Mathematics 71, Cambridge University Press, 1976.
\bibitem{Sharpe} D.~W.~Sharpe, \emph{Grade and the theory of linear
  equations}, Linear Algebra Appl.~18 (1977), 25--32.
\end{thebibliography}
\end{document}
